\clearpage\begingroup\NoSourceHead\SourceLeaf{499b}
\renewcommand{\thefootnote}{\arabic{footnote})}\setcounter{footnote}{58}
\ArithmeticHead{42}{L}{BOLZA, GAUSS UND DIE VARIATIONSRECHNUNG.}
\begin{center}\textls[100]{III. Abschnitt: Wirkungen der \textsc{Gauss}schen Arbeit.}\end{center}
\noindent\hspace*{\parindent}22. Bald nach der Veröffentlichung der \textsc{Gauss}schen Arbeit (1830)\footnote{Am 18. April 1830 schreibt \textsc{Gauss} an \textsc{Schumacher}: »Hierneben erhalten Sie einen der soeben
mir zugegangenen Abdrücke meiner Principia generalia . . .«, \textit{Briefwechsel zwischen C. F. Gauss und H. C.
Schumacher}, Bd. II, S. 230.}
setzten eine Reihe von umfangreichen und wichtigen Arbeiten über die allgemeine
Theorie der Extrema von Doppel- und mehrfachen Integralen ein,
die wir zum grössten Teil bereits im Vorhergehenden kurz erwähnt haben:
\textsc{Bordoni} (1831), \textsc{Poisson} (1833; datiert 10. November 1831), \textsc{Pagani} (1835;
datiert 15. Dezember 1834), \textsc{Ostrogradsky} (1836; datiert 24. Januar 1834),
\textsc{Delaunay} (1843; eingereicht vor 1. April 1842), \textsc{Cauchy}\footnote{\textit{Mémoires sur le calcul des variations}, Exercices d’analyse et de physique mathématique, T. III
(1844), S. 50—130.} (1844), \textsc{Sarrus}
(1848; eingereicht vor 1. April 1842), \textsc{Lindelöf}\footnote{\textit{Variationskalkylens theori och dess användning}, Helsingfors 1856.} (1856).

Da drängt sich von selbst die Vermutung auf, dass die \textsc{Gauss}sche Arbeit
den Anstoss zu dieser ganzen Entwicklungsreihe gegeben hat. Demgegenüber
ist jedoch zunächst festzustellen, dass keine dieser Arbeiten \textsc{Gauss} erwähnt,
mit Ausnahme derjenigen von \textsc{Pagani}. Bei der damals üblichen Laxheit im
Zitieren von Vorgängern beweist das allerdings nicht viel.

Im einzelnen ist zu sagen, dass bei \textsc{Bordoni} innere Gründe, die wir bereits
oben (Nr. 16) angeführt haben, gegen die Annahme sprechen, dass er die
\textsc{Gauss}sche Arbeit gekannt habe. Auch dürfte die Abfassungszeit seiner Arbeit
\clearpage\SourceLeaf{500}
\ArithmeticHead{43}{R}{WIRKUNGEN DER GAUSSSCHEN ARBEIT.}
\noindent kaum viel später als die Veröffentlichung der \textsc{Gauss}schen Arbeit fallen.
Anders liegen die Dinge bei \textsc{Poisson}. Zwischen dem Erscheinen der \textsc{Gauss}schen
Arbeit und der Datierung der \textsc{Poisson}schen liegen mehr als anderthalb
Jahre. Da wäre es immerhin möglich, dass \textsc{Poisson}, der sich schon 1816
erfolgreich mit der Variation von Doppelintegralen beschäftigt hatte, ohne
jedoch über die Schwierigkeiten der partiellen Integration hinauszukommen,
bei \textsc{Gauss} die Lösung dieser Schwierigkeiten gefunden hätte und dadurch zur
erneuten Inangriffnahme der Theorie der Extrema von Doppelintegralen veranlasst
worden wäre. Eine gewisse Bestätigung erhält diese Vermutung dadurch,
dass \textsc{Poisson} in der Vorrede zu seinem im Jahr 1831 erschienenen Buch
\textit{Nouvelle théorie de l’action capillaire} die Arbeit von \textsc{Gauss} zitiert\footnote{loc. cit. S. 7: »Depuis que cet ouvrage est écrit, j’ai eu connaissance d’un Mémoire de M. \textsc{Gauss},
qui paraît en ce moment sous le titre de Principia generalia theoriae figurae fluidorum in statu aequilibrii
(Gottingue 1830)«. \textsc{Poisson} gibt dann eine gedrängte Inhaltsangabe der \textsc{Gauss}schen Arbeit, die mit den
Worten schliesst: »Par les règles connues du calcul des variations, on détermine la surface inconnue du
liquide qui rend cette somme [nämlich die von \textsc{Gauss} mit $W$ bezeichnete Summe] un minimum, et, comme
on sait, on trouve à la fois l’équation générale de cette surface et l’équation particulière de son contour,
ce qui est l’avantage caractéristique de la méthode que M. \textsc{Gauss} a suivie«. Hier sind die Worte »comme
on sait« unverständlich, da man ja vor \textsc{Gauss} gerade \textls[100]{nicht} wusste, wie man die Randbedingungen erhalten
sollte. Schon \textsc{Todhunter} hat seine Verwunderung darüber ausgesprochen, dass \textsc{Poisson} in seiner Variationsarbeit
\textsc{Gauss} nicht erwähnt, wohl aber in der Kapillaritätsarbeit, vgl. \textsc{Todhunter}, loc. cit. 5), S. 54.}. Aber so
lange das Datum dieser Vorrede und der Zeitpunkt, in dem \textsc{Poisson} seine
grosse Arbeit über Variationsrechnung begonnen hat, unbekannt sind (und
diese Daten werden jetzt wohl kaum mehr festzustellen sein), muss die Frage
nach der Richtigkeit der obigen Vermutung unentschieden bleiben.

\textsc{Pagani} knüpft direkt an die Arbeit von \textsc{Gauss} an. Er stellt sich die
Aufgabe, die \textsc{Gauss}schen Resultate aus der allgemeinen Theorie der Variation
der mehrfachen Integrale herzuleiten, eine Theorie, die er in mehreren Punkten
über \textsc{Poisson} hinaus weitergeführt hat. Er verallgemeinert zunächst die \textsc{Lagrange}sche
Fundamentalformel (21) auf mehrfache Integrale, gibt dann weiter
für zweifache und dreifache Integrale die Umformung der ersten Variation
mittels partieller Integration auf Grund der \textsc{Gauss}schen Formeln (41) und (42)
und wendet schliesslich seine allgemeinen Resultate auf das \textsc{Gauss}sche Problem
an, worüber wir noch weiter unten berichten werden. Hier liegt also eine
unzweifelhafte, direkte Einwirkung von \textsc{Gauss} vor. Leider ist aber die Arbeit
von \textsc{Pagani} in der Geschichte der Variationsrechnung kaum beachtet worden,
\par\smallskip\noindent\hfill{\small 6\textsuperscript{*}}\par
\clearpage\SourceLeaf{501}
\ArithmeticHead{44}{L}{BOLZA, GAUSS UND DIE VARIATIONSRECHNUNG.}
\noindent obgleich sie einen beträchtlichen Teil der Resultate von \textsc{Ostrogradsky} der
Zeit der Veröffentlichung nach antizipiert hat. Im Gegensatz dazu zeigen die
Arbeit von \textsc{Ostrogradsky} und die übrigen oben angeführten Arbeiten keine
direkte Einwirkung von \textsc{Gauss}; wohl aber sind sie ohne Zweifel direkt oder
indirekt durch die Arbeit von \textsc{Poisson} angeregt worden.

Nach alledem hängt die Frage nach dem Umfang der Wirkung der \textsc{Gauss}schen
Arbeit auf die allgemeine Theorie der Extrema von mehrfachen Integralen
ganz davon ab, ob die \textsc{Poisson}sche Arbeit durch \textsc{Gauss} beeinflusst
worden ist oder nicht. Im ersteren Fall hat \textsc{Gauss} tatsächlich einen nachhaltigen
Einfluss ausgeübt, indem er durch Vermittlung von \textsc{Poisson} den Anstoss
zu der ganzen Reihe der oben genannten Arbeiten über Extrema von
Doppelintegralen gegeben hat. Im entgegengesetzten Fall muss man aber
sagen, dass die \textsc{Gauss}sche Arbeit fast ohne allen Einfluss auf die Entwicklung
der allgemeinen Theorie geblieben ist. Das muss man dann wohl dem Umstand
zuschreiben, dass \textsc{Gauss} selbst die Arbeiten seiner Vorgänger ganz
unberücksichtigt gelassen hat, und dass seine Methoden allzusehr auf das
spezielle in der Kapillaritätstheorie auftretende Variationsproblem zugeschnitten
waren.

23. Die vorangegangenen Ausführungen bezogen sich auf die Wirkung
der \textsc{Gauss}schen Arbeit auf die \textls[100]{allgemeine} Theorie der Extrema von Doppelintegralen.
Es bleibt uns nun noch übrig, über diejenigen Arbeiten zu berichten,
die sich mit dem \textls[100]{speziellen} \textsc{Gauss}schen Variationsproblem beschäftigen.
Es handelt sich bei denselben im wesentlichen teils um Vereinfachungen,
teils um Verallgemeinerungen der \textsc{Gauss}schen Beweise.

Hier sind zunächst einige Arbeiten zu nennen, in denen die allgemeine
Theorie der Variation der Doppelintegrale auf das \textsc{Gauss}sche Variationsproblem
angewandt wird, so \textsc{Pagani}\footnote{\textsc{Pagani}, loc. cit. 26), S. 96—99.} in der schon mehrfach erwähnten Arbeit und \textsc{Mainardi}
(1851) in einem Abschnitt einer grösseren Arbeit \textit{Ricerche sul calcolo delle
variazioni}\footnote{Annali di scienze matematiche e fisiche (\textsc{Tortolini}), Bd. III (1852), S. 187—190. \textsc{Mainardi}
macht dabei in der \textsc{Green}schen Formel denselben Vorzeichenfehler wie \textsc{Bordoni} (Nr. 16) und erhält das
richtige Endresultat nur dadurch, dass er diesen ersten Fehler durch einen zweiten in der Gleichung (59)
(oben, S. 40) kompensiert.}. Beide machen dabei keinen Gebrauch von der fertigen Formel (58)
\clearpage\SourceLeaf{502}
\ArithmeticHead{45}{R}{WIRKUNGEN DER GAUSSSCHEN ARBEIT.}
\noindent von \textsc{Bordoni} und \textsc{Poisson}, sondern leiten die Randbedingung durch besondere
Überlegungen her.

Demgegenüber erhält \textsc{Jellett}\footnote{loc. cit. 13a), S. 280.} (1850) die Randbedingung zwar nicht
für das \textsc{Gauss}sche Problem, das er merkwürdiger Weise nicht erwähnt, wohl
aber für das allgemeinere Problem des Extremums des Doppelintegrals
\[
\iint[\mu(x,y,z)\sqrt{1+p^2+q^2}+\mu'(x,y,z)]\,dx\,dy
\]
unmittelbar durch Anwendung der Formel (58), wobei sich die Randbedingung
mit einem Minimum von Rechnung in der Form
\[
\frac{1+pg_x+qg_y}{\sqrt{1+p^2+q^2}\sqrt{1+g_x^2+g_y^2}}=\text{Const.}
\]
ergibt.
\par\clearpage\endgroup
